% !PRINTED 190
\weberpage{190}

For the factors of $f(x, y)$, which we shall now encounter more often, we introduce as an abbreviation the sign

\weq{6}{x\beta_1 - y\alpha_1 = (x\beta_1).}

Likewise we shall set

\weq{7}{x_1 y_2 - x_2 y_1 = (x_1 y_2),}

in which $x_1, y_1$ and $x_2, y_2$ are two arbitrary pairs of variables. When we employ the inhomogeneous representation, we set still more briefly

\weq{8}{\begin{aligned}
x - \alpha_1 &= (0, 1),\ x - \alpha_2 = (0, 2) \ldots x - \alpha_n = (0, n) \\
\alpha_1 - \alpha_2 &= (1, 2) \ldots \alpha_i - \alpha_k = (i, k).
\end{aligned}}

If we make in the form (2) a linear substitution

\weq{9}{\begin{aligned}
x &= \alpha x' + \beta y' \\
y &= \gamma x' + \delta y'
\end{aligned}}

with the determinant, different from zero,

\weq{10}{r = \alpha\delta - \beta\gamma,}

then the form $f(x, y)$ passes over into a form $F(x', y')$ of degree $n$, and we have

\weq{11}{F(x', y') = a'_0 x'^n + a'_1 x'^{n-1} y' + a'_2 x'^{n-2} y'^2 + \cdots a'_n y'^n.}

The coefficients $a'_0, a'_1 \ldots a'_n$ of the transformed form are obtained as linear and homogeneous expressions in the coefficients $a_0, a_1 \ldots a_n$, and as homogeneous expressions of degree $n$ in the substitution coefficients $\alpha, \beta, \gamma, \delta$. One thus obtains, on setting $x' = 1, y' = 0$ or $x' = 0, y' = 1$,

\weq{12}{a'_0 = f(\alpha, \gamma), \qquad a'_n = f(\beta, \delta).}

The remaining coefficients $a'$ can be expressed by means of the derivatives of the function $f(x, y)$, e.g.

% [sic] The print has the factor (n - 1); the coefficient form in (2) does not support it algebraically.
\[
(n - 1) a'_1 = \beta f'(\alpha) + \delta f'(\gamma),
\]

which we shall not carry out further.

If we apply the transformation (9) to the two pairs of variables $x, y$ and $\alpha_k, \beta_k$ simultaneously, we obtain the linear factors of $f(x, y)$ transformed, namely

\weq{13}{(x'\beta'_k) = r(x\beta_k),}

as one finds from the multiplication rule for determinants, or also by direct computation. These factors $(x\beta_k)$ are therefore also \gesperrt{covariants} of $f(x, y)$, though to be sure irrational ones, since they do not depend rationally on the coefficients of $f(x)$. In the same way one obtains
%CONT>>
